\section{Background and Related Work}\label{sec:related}

\subsection{What the First Bytes Say}\label{sec:signatures}

A listener can decide a protocol from the first bytes only if the client sends first. For the
protocols this paper serves directly, the client does.
\begin{itemize}
\item \textbf{TLS.} The client opens with a handshake record that carries a
ClientHello~\cite{rescorla2018rfc8446}. The record's content type and the handshake message type
are fixed. The record's version field is a legacy field that receivers must ignore, so only its
major version is a usable key.
\item \textbf{HTTP/2 with prior knowledge.} The client opens with a fixed connection preface of
\HtwocNeed{} bytes~\cite{thomson2022rfc9113}. The other way to start HTTP/2 on cleartext, the
\texttt{Upgrade} header, is deprecated.
\item \textbf{HTTP/1.1.} The client opens with a request line: a method token, one space, a
target and a version~\cite{fielding2022rfc9112}. A server should ignore at least one empty line
before it. The methods are those of HTTP semantics~\cite{fielding2022rfc9110} and
PATCH~\cite{dusseault2010rfc5789}.
\item \textbf{SSH.} Each side sends an identification string that starts with
\texttt{SSH-}~\cite{ylonen2006rfc4253}. The standard does not say which side sends first, so a
first-bytes listener can recognise only clients that send first.
\item \textbf{MQTT.} The first packet of a client must be CONNECT. Its protocol name follows a
Remaining Length field of variable size, so the name's offset depends on the
length~\cite{banks2014mqtt311,banks2019mqtt5}.
\item \textbf{PROXY.} A proxy may prepend a header that carries the original addresses, in a text
form (v1) or a binary form (v2). The specification requires the receiver to be
configured to expect it, and never to guess whether it is present~\cite{tarreau2026proxy}.
\end{itemize}
Server-speaks-first protocols, such as SMTP, have the server send a greeting
first~\cite{klensin2008rfc5321}. A compliant client of such a protocol waits. Its connection
delivers no byte from which to decide, so a first-bytes listener can serve it only after a timer.

The signatures overlap in ways a detector must handle. The binary PROXY header starts with an
empty line, which an HTTP/1.1 parser would skip. The HTTP/2 preface parses as an HTTP/1 request
line with the method \texttt{PRI}. \texttt{SSH-} is a legal prefix of an HTTP method token. An
inline command of the Redis protocol, such as \texttt{GET key}, shares its first bytes with an
HTTP request. A server that tells such classes apart needs a rule that does not depend on the
order in which it tries them.

\subsection{Systems That Detect a Protocol from Its First Bytes}\label{sec:systems}

We read the documentation and the source code of the systems below at the versions cited. Each
statement here holds for that version.

\runin{Proxies.} nginx's stream module reads the ClientHello before any handshake to route by its
server name~\cite{nginx2026}. It peeks with \texttt{MSG\_PEEK} only when it does not terminate TLS
and its events are edge-triggered, on epoll with peer half-close reporting or on kqueue. Otherwise
it reads the bytes and replays them to the upstream. HAProxy inspects content in its
TCP mode by reading into its buffer and evaluating its rules again on each new chunk, for at most
its inspection delay~\cite{haproxy2026}. A connection that matches no rule passes to the default
back end; the documentation discusses that delay with SMTP as its example. Envoy runs listener
filters on the accepted socket~\cite{envoy2026}. Their shared buffer peeks with
\texttt{MSG\_PEEK}; of the filters we read, only the PROXY filter consumes bytes. caddy-l4 tries ordered routes of matchers on a recorded
buffer that it rewinds after each matcher, on top of Caddy~\cite{caddyl42026,caddy2026}. sslh reads
the client's bytes, runs its probes on everything read so far, and connects to a configured
protocol when its timeout expires~\cite{sslh2026}. sshttp peeks once and chooses between SSH and
HTTP or HTTPS~\cite{sshttp2023}. Traefik splits TLS from other traffic by its first byte and reads
and replays the bytes~\cite{traefik2026}.

\runin{Libraries.} Netty's port unification example reads a few bytes without consuming them,
rewires the connection's pipeline, and hands the accumulated bytes to the next
handler~\cite{netty2026}. Jetty's detector asks each detecting connection factory in turn and
copies the unread bytes into a new buffer for the chosen one~\cite{jetty2026}. cmux records the
bytes it reads and lets each matcher start again from the first byte~\cite{cmux2021}.
hyper-util's automatic connection reads at most the HTTP/2 preface to choose between HTTP/1 and
HTTP/2, and replays it~\cite{hyperutil2026}. Go's HTTP server checks the preface in user space
when unencrypted HTTP/2 is enabled~\cite{go2026}.

\runin{What they share.} Most of these systems read the first bytes into user space and replay
them; Envoy, and nginx on some paths, peek instead. None that we read documents a comparison of
the two. A server-speaks-first protocol shares a detecting listener in two ways only. One is a
timer, as HAProxy's default back end, sslh's timeout target and Envoy's default filter chain
provide. The other is a listener that speaks first itself, as sshttp's SMTP mode does. Built-in detectors for SSH or MQTT
appear in proxies (sslh and caddy-l4 for SSH, HAProxy for MQTT) and in none of the libraries
above. Their documents do not say what happens when the first read is shorter than a signature.
In the code, most of them wait for more bytes; sshttp decides on its single peek.

\subsection{Kernel and Operating System Mechanisms}\label{sec:kernel}

On Linux, a group of sockets can share a port with \texttt{SO\_REUSEPORT}. A program attached to
the group selects the listener for a new TCP connection at its SYN, before any application byte
exists~\cite{man2026socket,linux2026v72}. A socket-lookup program selects a socket by
addresses, ports and interface, and sees no payload~\cite{kernel2026sklookup}.
\texttt{TCP\_DEFER\_ACCEPT} delays the accept until data arrives, on the same
socket~\cite{man2026tcp}. A socket map can redirect the data of an established socket, not the
connection~\cite{kernel2026sockmap}. FreeBSD's accept filters hold a connection back from
\texttt{accept} until a condition on its data holds~\cite{freebsd2026acceptfilter}. On Windows,
\texttt{AcceptEx} can return the first block of data with the accepted socket; with a receive
buffer, it completes only after data arrives~\cite{microsoft2026acceptex}. In the documentation
and source code we read, none of these mechanisms routes a TCP connection to a listener by its
first application bytes. They decide earlier, or they only delay.

\subsection{Standards and Research}\label{sec:research}

TCPMUX served many services on one TCP port by having the client name the service in its first
line~\cite{lottor1988rfc1078}. It is now historic~\cite{zimmermann2016rfc7805}. ALPN negotiates
the application protocol inside the TLS handshake, for several protocols behind one
port~\cite{friedl2014rfc7301}. For UDP, a receiver can demultiplex several protocols by the first
byte of each datagram~\cite{aboba2023rfc9443}.

We found no peer-reviewed paper whose subject is a server that serves several protocols on one
listener. Four lines of work touch the question. ALPACA shows that TLS does not bind a connection
to the application protocol its client intended~\cite{brinkmann2021alpaca}. Work on
probe-resistant proxies shows that how a server reacts to first bytes it does not recognise,
including its timeouts, is a fingerprint~\cite{frolov2020probe,frolov2020httpt}. An analysis of
STARTTLS finds defects where bytes read before a protocol switch are processed after
it~\cite{poddebniak2021starttls}. A first-bytes server hands bytes it has already read to the next
protocol, and must do so correctly. Traffic classification infers an application protocol from
the content or the first packets of a flow, on monitors rather than on
servers~\cite{dreger2006dynamic,deri2014ndpi,bernaille2006fly}.

\subsection{Gaps}\label{sec:gaps}

Four gaps follow from the sources above. No system we read publishes a measured cost of its
detection. None documents a comparison of peek with read and replay. Server-speaks-first
protocols are served only through a timer, or by a listener that speaks first. What happens when
the first read is shorter than a signature is mostly undocumented. This paper measures the first
two for one server. It turns the last two into rules of that server, stated before any run and
tested on a table of hard cases.
